\documentclass[a4paper,12pt]{article}

% Pakete für zusätzliche Funktionen
\usepackage[utf8]{inputenc}   % UTF-8 Zeichencodierung
\usepackage[T1]{fontenc}     % Korrekte Silbentrennung und Umlaute
\usepackage[ngerman]{babel}  % Deutsche Spracheinstellungen
\usepackage{amsmath, amssymb} % Mathematische Symbole
\usepackage{graphicx}        % Einfügen von Bildern
\usepackage{hyperref}        % Hyperlinks
\usepackage{geometry}        % Seitenränder
\geometry{a4paper, margin=2.5cm}
\usepackage{listings}        % Code-Listings
\usepackage{csquotes}        % Zitate
\usepackage{xurl}
% Einstellungen für Hyperlinks
\hypersetup{
    colorlinks=true,
    linkcolor=blue,
    citecolor=blue,
    filecolor=magenta,
    urlcolor=cyan,
    pdftitle={Webtechnologies and Applications (WTA)},
    pdfpagemode=FullScreen,
}

% Titel und Autoreninformationen
\title{\textbf{Webtechnologies and Applications (WTA)}}
\author{Kürsat Darcan | MFWS422A}
\date{Abgabedatum: \today}

\begin{document}

% Deckblatt
\maketitle
\thispagestyle{empty}
\vspace{2cm}
\begin{center}
    \includegraphics[width=0.3\textwidth]{FHDW_Logo_RGB-01.svg.png} % Ersetzen Sie "example-image" durch Ihr Logo/Bild
    \\
    \vspace{1cm}
    \textbf{Studiengang: Wirtschaftsinformatik}\\
    \textbf{Fachhochschule der Wirtschaft (FHDW)}
\end{center}
\newpage

% Inhaltsverzeichnis
\renewcommand{\thepage}{\roman{page}} % Seitenzahlen in römischen Ziffern für Verzeichnisse
\tableofcontents
\newpage

% Bei Bedarf Abbildungsverzeichnis
%\listoffigures
%\addcontentsline{toc}{section}{Abbildungsverzeichnis}
%\newpage

% Bei Bedarf Tabellenverzeichnis
%\listoftables
%\addcontentsline{toc}{section}{Tabellenverzeichnis}
%\newpage

% Abkürzungsverzeichnis
%\section*{Abkürzungsverzeichnis}
%\addcontentsline{toc}{section}{Abkürzungsverzeichnis}
%\begin{description}
%    \item[HTML] Hypertext Markup Language
%\end{description}
\newpage

% Hauptteil mit arabischen Seitenzahlen
\renewcommand{\thepage}{\arabic{page}} % Seitenzahlen in arabischen Ziffern für den Hauptteil
\setcounter{page}{1}


% Einleitung
\section{Einleitung}
\subsection{Hintergrund und Relevanz des Themas}
Mit dem Fortschritt der künstlichen Intelligenz (KI) entwickeln sich auch Deepfake-Technologien stetig weiter, 
wodurch die öffentliche Aufmerksamkeit für diese Thematik zunimmt. 
Diese Technologien ermöglichen die Manipulation digitaler Medieninhalte wie Videos, 
Bilder und Audiodateien, sodass diese täuschend echt wirken, 
obwohl sie nie in dieser Form erstellt wurden. 
Deepfakes sind jedoch keine eigenständige Technologie, 
sondern basieren auf einer Kombination verschiedener KI-Methoden. \cite{BVDW2024}\\
\subsection{Problemstellung und Herausfoderung}
Die vielseitigen Einsatzmöglichkeiten von Deepfakes reichen von kreativen Anwendungsbereichen, 
wie der Erstellung von Unterhaltung, bis hin zu manipulativen oder schädlichen Absichten. 
So können Deepfakes gezielt eingesetzt werden, um die öffentliche Meinung zu beeinflussen, 
Personen oder Unternehmen zu schädigen oder Desinformation zu verbreiten. \cite{CounterExtremism2020}\\
\subsection{Ziel der Ausarbeitung und Aufbau}
Diese Ausarbeitung untersucht mögliche Gegenmaßnahmen zur Erkennung und Bekämpfung von Deepfake-Technologien. 
Dabei wird auch die ethische Seite betrachtet, 
insbesondere die Frage, welche Verantwortung Entwickler und Nutzer dieser Technologie tragen. 
Abschließend wird ein Ausblick auf die zukünftige Entwicklung der Technologie und deren mögliche Auswirkungen auf die Gesellschaft gegeben.
\newpage

% Gegenmaßnahmen gegen Deepfakes
\section{Gegenmaßnahmen gegen Deepfakes}
Die Verbreitung von Deepfakes stellt eine erhebliche Herausforderung dar, 
da sie unterschiedliche Auswirkungen auf die Gesellschaft haben kann. 
Neben der Gefahr von Desinformation und Betrug kann auch das Vertrauen in digitale Medien nachhaltig geschädigt werden. 
Daher ist es essenziell, Methoden zur Erkennung von Deepfakes zu entwickeln und Maßnahmen zur Regulierung einzuführen. \cite{CounterExtremism2020}

\subsection{KI-Erkennungsmechanismen}
Mit dem Fortschritt der KI-Technologie werden spezialisierte Modelle zur Erkennung von Deepfakes entwickelt wie zum beispiel
Convolutional Neural Networks (CNNs).
Die CNNs wird spezial auf 
Anomalien in Gesichtsbewegungen, 
Lichtreflexionen oder Sprachmodulationen traniert, damit  
die auf Fälschungen hinweisen könnten. 
Darüber hinaus werden auch Transformer-Modelle und hybride Ansätze genutzt, 
um Deepfakes zuverlässiger zu enttarnen. \cite{BVDW2024}

\subsection{Deklarierung von KI-generierten Medien}


\subsection{Prävention und Bildung}
\subsubsection{Öffentliche Aufklärung}
\newpage

% Ethik und Deepfakes
\section{moralische Herausforderungen}
\subsection{Verantwortung der Entwickler}
\subsubsection{Innovation vs. Missbrauch}
\subsubsection{Regulierung von Open-Source-Projekten}
\subsection{Gesellschaftliche Folgen}
\subsubsection{Manipulation der öffentlichen Meinung}
\subsubsection{Persönliche Schäden}
\newpage

% Zukunftsperspektiven
\section{Zukunftsperspektiven}s
\subsection{Technologische Entwicklungen}
\subsection{Gesellschaftliche Anpassungen}
\newpage

%Fazit
\section{Fazit}



% Literaturverzeichnis
\newpage
\addcontentsline{toc}{section}{Literaturverzeichnis}
\section*{Literaturverzeichnis}
\begin{thebibliography}{99}
    \bibitem{Beispiel1} Beispielautor, B. (2025). 
    \textit{Titel des Buches}. Verlag, Ort.
    \bibitem{Beispiel2} Beispielautor, C. und Beispielautor, D. (2024). Artikelname. 
    \textit{Zeitschrift}, 12(3), 45-67.
    \bibitem{BSI2025} Bundesamt f\"ur Sicherheit in der Informationstechnik (BSI). (2025). Deepfakes – Bedrohungspotenziale und Gegenmaßnahmen. 
    Verfügbar unter: \url{https://www.bsi.bund.de/DE/Themen/Unternehmen-und-Organisationen/Informationen-und-Empfehlungen/Kuenstliche-Intelligenz/Deepfakes/deepfakes_node.html} (zuletzt aufgerufen am 21.01.2025). Kopie aufbewahrt als \texttt{BSI - Deepfakes - Gefahren und Gegenmaßnahmen.pdf}.
    \bibitem{BVDW2024} Bundesverband Digitale Wirtschaft (BVDW). (2024). Deepfakes – Eine juristische Einordnung. 
    Verfügbar unter: \url{https://www.bvdw.org/wp-content/uploads/2024/07/2024_BVDW_Deepfakes.pdf} (zuletzt aufgerufen am 22.01.2025). Kopie aufbewahrt als \texttt{2024\_BVDW\_Deepfakes\_eine\_einordnung.pdf}.
    \bibitem{CounterExtremism2020} Counter Extremism Project. (2020). Deepfakes – Bedrohungspotenziale und Gegenmaßnahmen.
    Verfügbar unter: \url{https://www.counterextremism.com/sites/default/files/200806_Deep_Fakes_DE_WEB_DS.pdf} (zuletzt aufgerufen am 22.01.2025). Kopie aufbewahrt als \texttt{200806\_Deep\_Fakes\_DE\_WEB\_DS.pdf}.
    \bibitem{Kaur2024} Kaur, A., Noori Hoshyar, A., Saikrishna, V., Firmin, S., \& Xia, F. (2024). Deepfake video detection: challenges and opportunities. Verfügbar unter: 
    \url{https://link.springer.com/content/pdf/10.1007/s10462-024-10810-6.pdf} (zuletzt aufgerufen am 23.01.2025). Kopie aufbewahrt als \texttt{Deepfake\_video\_detection\_challenges\_and\_opportunities.pdf}.
    \bibitem{Arxiv2020} Nguyen, H. H., Yamagishi, J., \&< Echizen, I. (2020). Capsule-Forensics: Using Capsule Networks to Detect Forged Images and Videos. 
    Verfügbar unter: \url{https://arxiv.org/pdf/2003.08685} (zuletzt aufgerufen am 23.01.2025). 
    \bibitem{BPB2024} Bundeszentrale für politische Bildung (BPB). (2024). Regulierung von Deepfakes. 
    Verfügbar unter: \url{https://www.bpb.de/lernen/bewegtbild-und-politische-bildung/556822/regulierung-von-deepfakes/} (zuletzt aufgerufen am 23.01.2025). 
    \bibitem{Springer2024} Dittmann, J., \& Vielhauer, C. (2024). Technische Ansätze zur Deepfake-Erkennung. In: Künstliche Intelligenz und digitale Ethik. Springer. 
    Verfügbar unter: \url{https://link.springer.com/chapter/10.1007/978-3-662-65964-9_10} (zuletzt aufgerufen am 23.01.2025).
    \bibitem{Fraunhofer2024} Fraunhofer ISI. (2024). Chancen und Risiken von Deepfakes für Politik, Wirtschaft und Gesellschaft. 
    Verfügbar unter: \url{https://www.isi.fraunhofer.de/de/presse/2024/presseinfo-17-deepfakes-chancen-risiken-Politik-Wirtschaft-Gesellschaft.html} (zuletzt aufgerufen am 23.01.2025). 
    \bibitem{Tibor2024} Tibor.net. (2024). Instagram markiert Posts mit „KI-generiert“ – was bedeutet das? 
    Verfügbar unter: \url{https://tibor.net/digital/kuenstliche-intelligenz/instagram-markiert-posts-mit-mit-ki-generiert-was-bedeutet-das/} (zuletzt aufgerufen am 23.01.2025).
    \bibitem{Bundesdruckerei2024} Bundesdruckerei. (2024). Desinformation und Fake-ID. 
    Verfügbar unter: \url{https://www.bundesdruckerei.de/de/innovation-hub/desinformation-und-fake-id} (zuletzt aufgerufen am 23.01.2025). 
    \bibitem{BPBDeepfake2024} Bundeszentrale für politische Bildung (BPB). (2024). Technische Ansätze zur Deepfake-Erkennung und -Prävention. 
    Verfügbar unter: \url{https://www.bpb.de/lernen/bewegtbild-und-politische-bildung/556855/technische-ansaetze-zur-deepfake-erkennung-und-praevention/} (zuletzt aufgerufen am 23.01.2025).
\end{thebibliography}


\newpage
% Anhang
\appendix
\section{Anhang}
Hier können zusätzliche Informationen, wie Code-Beispiele oder ausführliche Tabellen, eingefügt werden.

% Ehrenwörtliche Erklärung
\newpage
\addcontentsline{toc}{section}{Ehrenwörtliche Erklärung}
\section*{\texttt{Ehrenwörtliche Erklärung}}
Hiermit erkläre ich, dass ich die vorliegende schriftliche Ausarbeitung im Modul \textbf{Webtechnologies and Applications} selbstständig
angefertigt habe. Es wurden nur die in der Arbeit ausdrücklich benannten Quellen und
Hilfsmittel benutzt. Wörtlich oder sinngemäß übernommenes Gedankengut habe ich als
solches kenntlich gemacht. Diese Arbeit hat in gleicher oder ähnlicher Form noch keiner
Prüfungsbehörde vorgelegen.

\vspace{3cm}
\noindent\begin{tabular}{p{0.5\textwidth}p{0.5\textwidth}}
    \hrulefill & \hrulefill \\
    Ort, Datum & Unterschrift \\
\end{tabular}

\end{document}
